\section*{The Propers: Detailed Commentary}

\subsection*{The typical edition against the Benziger printing}

The Vatican \latin{editio typica} controls every word printed above. Against the
Benziger \latin{editio iuxta typicam} of the same year, whose text reaches this
guide only in an uncorrected optical layer, and allowing for that layer's accent
substitutions, the two agree through the
Introit, the Collect, the body of the Epistle, the Alleluia, the Offertory, the
body of the Secret and the Communion --- including \latin{ad me} and
\latin{miserére mihi}, the semicolon after \latin{consístere}, \latin{super omnem
terram}, \latin{respéxit me}, \latin{hymnum}, \latin{dédero} and \latin{sǽculi}.

\textbf{Where they differ, one difference is a word and the rest are pointing,
abbreviation and cue length.} At Lk. 7:14 the typical edition prints
\latin{Aduléscens} with \emph{u} and Benziger \latin{Adolescens} with \emph{o};
the Clementine reads the \emph{o}, and so do the uncorrected optical layers of
the Venice 1570 and the 1604 \latin{typica}, so the \emph{u} is the 1962 typical
edition's own form. At the same verse
the stop stands outside the closing parenthesis here, \latin{stetérunt).}, and
inside it there, where the Clementine agrees with Benziger. Benziger adds two
commas at Gal. 6:6 and one after \latin{Mentes nostras}, prints the
antiphon-repetition cue \latin{Inclína} after the Introit's \latin{Glória
Patri.}, and abbreviates both the Secret's and the Postcommunion's conclusions to
\latin{Per Dóminum.} where the typical edition writes them out to \latin{in
unitáte.} --- the reverse of the Fourteenth Sunday's direction, where it is
Benziger that prints the longer cue.

\subsubsection*{Where the appointed text is not the Bible's text}

The table prints departures at six elements. Nine of them are word-level:
\latin{ad me} and \latin{miserére mihi} at the Introit, with \latin{quia} for
\latin{quoniam} in its psalm verse; \latin{super omnem terram} at the Alleluia;
\latin{respéxit me}, \latin{deprecatiónem meam} and \latin{hymnum} at the
Offertory; \latin{dédero} and \latin{sǽculi} at the Communion. The eight outside
the psalm verse stand also in the uncorrected optical layers of the Pustet
Ratisbon 1862, the Venice 1570 and the Vatican \latin{typica} of 1604, so none
of the eight is a 1962 act. The rest of the table is matter of another kind ---
four extents cut, the lesson's prefixed address, the Gospel's substituted
incipit, capitalisation, and the stop that moves outside the parenthesis at Lk.
7:14 --- with one word form that runs the other way: the uncorrected optical
layers of the Venice 1570 and the 1604 \latin{typica} show \latin{Adolescens}
with the Clementine, and the \emph{u} of \latin{Aduléscens} is the 1962 typical
edition's own.

\begin{divergencetable}
Int., Ps. 85:1 & \latin{aurem tuam ad me, et exáudi me:} & \latin{aurem tuam et exaudi me,} --- the antiphon \textbf{adds} \latin{ad me}\\
Int., Ps. 85:1--2 & v.~1b and v.~2a absent; \latin{exáudi me:} runs into \latin{salvum fac} & \latin{quoniam inops et pauper sum ego. Custodi animam meam, quoniam sanctus sum;}\\
Int., Ps. 85:3 & \latin{miserére mihi, Dómine,} & \latin{Miserere mei, Domine,}\\
Int. verse, Ps. 85:4 & \latin{servi tui: quia ad te} & \latin{servi tui, quoniam ad te}\\
Ep., Gal. 5:25 & \latin{Fratres:} prefixed; \latin{spíritu} twice lower case & nothing corresponds to the address; \latin{Spiritu} capitalised twice\\
All., Ps. 94:3 & \textbf{\latin{super omnem terram.}} & \textbf{\latin{super omnes deos.}} --- a different predicate\\
Gosp., Lk. 7:11 & \latin{In illo témpore: Ibat Iesus in civitátem} & \latin{Et factum est: deinceps ibat in civitatem}\\
Gosp., Lk. 7:14 & \textbf{\latin{Aduléscens}}; \latin{stetérunt).} & \textbf{\latin{Adolescens}}; \latin{steterunt.)}\\
Off., Ps. 39:2 & \textbf{\latin{et respéxit me:}} & \textbf{\latin{et intendit mihi.}}\\
Off., Ps. 39:3 & \textbf{\latin{et exaudívit deprecatiónem meam:}} & \textbf{\latin{Et exaudivit preces meas,}}\\
Off., Ps. 39:3 & rest of v.~3 absent & \latin{et eduxit me de lacu miseriæ \dots\ et direxit gressus meos.} --- cut entire\\
Off., Ps. 39:4 & \textbf{\latin{hymnum Deo nostro.}} and stop & \textbf{\latin{carmen Deo nostro.}} followed by \latin{Videbunt multi \dots} --- v.~4b cut\\
Comm., Jn. 6:52 & v.~52a absent; \latin{Panis,} capitalised & \latin{Si quis manducaverit ex hoc pane, vivet in æternum: et panis}\\
Comm., Jn. 6:52 & \textbf{\latin{quem ego dédero}}; \textbf{\latin{pro sǽculi vita.}} & \textbf{\latin{quem ego dabo}}; \textbf{\latin{pro mundi vita.}}\\
\end{divergencetable}

Beside these stand three classes of difference that are not departures of
substance: pointing at the joins, where the Introit's colons stand at a cut that
spans a Clementine sentence end, \latin{exáudi me: salvum fac}, and at a
Clementine full stop, \latin{in te: miserére}, while the Gradual sets a full stop
and the versicle sign, \latin{Altíssime. ℣.~Ad annuntiándum}, where the
Clementine has a colon; orthography and word division, \latin{huiúsmodi}
for \latin{hujusmodi} and \latin{et si} for \latin{etsi} at Gal. 6:1; and
capitalisation, which in a modern electronic Clementine is that text's convention
and in the missal is that book's, so the difference is a convention of each book
and not a variant.
\textbf{The Epistle is remarkably clean:} after its liturgical address, the whole
of Gal. 5:25--6:10 matches the tracked Clementine word for word, the alignment
returning five difference runs across twelve verses --- two capitalisations, one
word division, one \emph{i}/\emph{j} orthography and one comma. The Gradual
departs only in pointing: the Clementine prints a comma after
\latin{misericórdiam tuam} and the 1962 typical edition does not, while the
uncorrected optical layers of the Benziger 1962, the Pustet 1862, the Venice
1570 and the 1604 \latin{typica} all show it.

\correction{Psalm numbering, and the rubrics behind the colour and the single oration}{%
The four psalms are cited in the missal's Vulgate numbering, and the tracked
Douay concordance supplies their English equivalents rather than arithmetic. Vulgate Pss.\ 85
and 94 answer English 86 and 95 with no offset, the Clementine counting the
inscription inside v.~1; Vulgate Pss.\ 39 and 91 answer English 40 and 92 at an
offset of one, their inscriptions standing as Vulgate v.~1 and going unnumbered
in English Bibles. Green is not printed at this Mass: it follows RG 127~b of the
1960 \latin{Rubricae generales}, whose September Ember exception is a rule about
\latin{feriæ} and does not reach a Sunday. That the formulary prints one Collect,
one Secret and one Postcommunion follows inside the 1960 code from RG 434~b with
RG 435; the instrument that abolished the seasonal second and third orations
themselves is not the \latin{Rubricae generales} but \work{Variationes in
Breviario et Missali Romano} cap.~IV n.~18, \latin{Orationes pro diversitate
Temporum abolentur}, at \work{AAS}~52 (1960) printed p.~709. The same
promulgation carries a \latin{Rubricae generales} n.~18 at printed p.~599, on
transferring the impeded Sundays after the Epiphany, so the bare numeral names two
different instruments.}

\subsection*{How this Mass came to be this Sunday's}

\subsubsection*{The three orations descend as a unit, and their wording has moved twice}

All three composed orations stand together, as one Mass and in liturgical order,
in the Gregorian Hadrianum as Wilson prints it: section XXXIIII, heading
\latin{DOMINICA .XVI. POST PENTECOSTEN.}, printed p.~174, with the manuscript
folio mark \latin{[fo. 168v]}. The same three stand in the Old Gelasian, Vat.
Reg. lat. 316, at Wilson's Book~III sect.~XI, printed p.~230, under the heading
\latin{Item alia Missa} and with no Sunday number at all. \textbf{Between the
eighth-century books and 1962 the wording moves in two places, and the two
ancient books do not move together}: both of them read \latin{contra diabolicos
tueantur semper incursus} where the 1962 Secret reads \latin{contra diabólicos
semper tueántur incúrsus}, and at the Postcommunion the Old Gelasian sets the
vocative differently, \latin{possideat, Domine, quaesumus}, where the Hadrianum
and 1962 read \latin{possídeat, quǽsumus, Dómine}. So the Hadrianum stands apart
from 1962 at one place and the Old Gelasian at two, and the two ancient witnesses
differ from each other at the second.

Two further variants are transmitted at the Postcommunion. Wilson's apparatus
records Pamelius reading \latin{iugiter aevi praeveniat} and the S(2) and Gerbert
witnesses \latin{iugiter eius proveniat} against the \latin{iugiter eius
praeveniat} of the received text, and notes that the Vatican manuscript's
\latin{sunt} for \latin{sed} is a scribal error corrected by Tommasi.

\subsubsection*{The Sunday number, and witnesses that disagree by one}

Five sacramentary witnesses number this Mass and they do not agree. Wilson's
apparatus at Gelasian III.XI records that ``This Missa is assigned by R.~S.\
Gerb.\ to the seventeenth Sunday after Pentecost: see the Missa for the sixteenth
Sunday in Pam., and that for the seventeenth Sunday in Men.,'' with the marginal
citations \latin{R. S. / Gerb. 175. / Pam. 411. / Men. 179.} --- \textbf{so Wilson
reports Rheinau, S.~Gallen, Gerbert and Menard at two Sundays later than 1962 and
Pamelius at one.} The page image settles what the optical layers lose: the second
Collect of the Gelasian Mass, \latin{Da, quaesumus, Domine, hanc mentem populo
tuo}, carries only \latin{R. S. / Gerb. 175.} in its margin, so Pamelius and
Menard do not have it, and the three prayers all five numbering witnesses carry
are exactly the three the 1962 formulary keeps. The same one-Sunday split
reproduces at the adjoining section, so the offset is a numeration base and not a
local accident.

The Liber Comitis heads the Mass \latin{Dominica mensis vi} at PL~30 cols.\
521--522, its count from the feast of St~Lawrence having stopped one Sunday
earlier, and Schuster calls this one \latin{Quinta post Sancti Laurentii}, the
last of the Lawrence Sundays; the two reckonings differ by one. What Migne prints
there under Jerome's name --- its date, its provenance, its textual authority,
and whether the \latin{mensis vi} and \latin{mensis viii} headings with no
\latin{mensis vii} between them are the book's own or Migne's arrangement --- is
unsettled, so the Epistle-and-Gospel pairing here is what the Liber Comitis
prints.

\subsubsection*{Honorius's Fifteenth Sunday, and the chants that did not move}

Honorius Augustodunensis, writing c.~1120--1130, has this Sunday's Collect where
1962 has it: at \work{Gemma animae} IV cap.~LXXVIII, \latin{Dominica decima
quinta, <Inclina:> sub lege}, PL~172 cols.\ 719--720, he writes \latin{cui Oratio,
Ecclesiam tuam, Domine, concordat, quam ut Deus mundet, muniat, rogat}, and
repeats the prayer at cap.~LXXIX. His Sixteenth offers \latin{Absolve, Domine,
tuorum vincula populorum, vel in Oratione, Tua nos, Domine}, so this Secret's
incipit stands there as an alternative oration. \textbf{His Alleluias are not
this one, and he says so himself}: \latin{Alleluia, Domine, exaudi \dots\ Alii
cantant Alleluia, Confitemini Domino \dots\ Alii cantant Alleluia, Paratum cor
meum}. The Alleluia 1962 sings here stands at his Fourteenth. Honorius's reading
of the whole Mass off the September historia of King Josiah is his own scheme
applied to a twelfth-century German use.

In the margins of Ottobonianus 313 the chant series and the oration series run
two Sundays apart across five consecutive sections, uniformly. The four-element
Mass standing under this Sunday's Introit in that margin --- Introit
\latin{Inclina domine aurem}, respond \latin{Bonum est confidere}, Offertory
\latin{Expectans expectaui}, Communion \latin{Qui manducat} --- agrees with
Honorius at three slots of four, and is not the 1962 four. Whether that margin's
\latin{Bonum est confidere} under this Introit is a respond shared across two
consecutive Sundays or a copying error is unsettled; the doubling on Wilson's
p.~173 is real, his own index printing ``173 (bis).'' Delisle assigns the
manuscript to the second half of the ninth century and Ebner follows him, though
Ebner's later judgment favoured an earlier date; Wilson reports both and adopts
neither.

The gregorien.info index of the six earliest Mass antiphonaries Hesbert collected
already places this Introit, Gradual and Offertory at \latin{Dominica XV post
Pentecosten}, AMS~187, but gives that Mass the Communion \latin{Qui manducat
carnem meam} and no fixed Alleluia at all; \latin{Panis quem ego dedero} stands
there at the Fourteenth Sunday. \textbf{So against that record the move runs one
way: the Communion this Mass sings stands at the Fourteenth Sunday in AMS, and
the Communion AMS gives the Fifteenth has not taken its place at the 1962
Fourteenth.} Within the tracked registry that Sunday sings \latin{Primum
quaerite} at Mt.\ 6:33, while the 1962 book's own \latin{Qui manducat meam
carnem}, Jn.\ 6:57, stands at the Ninth Sunday. Every antiphonary figure here is
the gregorien.info database's report of Hesbert's printed pages.

Guéranger's 1909 volume shows the formulary before the 1960 change, printing the
three seasonal orations by cross-reference --- ``The other Collects, as on
page~120'' at printed p.~345, ``The other Secrets, as on page~130'' and ``The
other Postcommunions, as on page~131'' at p.~355 --- and names the day ``the
Sunday of the widow of Naim.'' He prints at p.~356, under Vespers, the Magnificat
antiphon \latin{Propheta magnus surrexit in nobis, et quia Deus visitavit plebem
suam}; that antiphon is his book's. Schuster's three pages name no station church
for this Sunday.

\subsection*{Introit \cue{Int.}}

The antiphon is a three-segment centonisation: v.~1a as far as \latin{exáudi
me}, then v.~2b from \latin{salvum fac}, then v.~3 entire, with v.~4 as the psalm
verse. \textbf{The clause it drops first is the psalm's ascription}, and that
ascription is unique: the ascriptions of Vulgate Pss.\ 72--88 in the tracked
Clementine give Asaph at 72--82, \latin{filiis Core} at 83, 84, 86 and 87, Eman
the Ezrahite at 87 as well, Ethan the Ezrahite at 88, and
\latin{Oratio ipsi David.} at 85 alone. Ps.~85 is the single Davidic prayer in the
whole of Book~III, which the tracked text bounds by the doxologies at Ps.\
71:19--20 and Ps.\ 88:53. What the cento sings after that omission is a psalm
addressed by a servant to his God with no name attached to it.

\begin{witnessnote}
\wlocus{Augustine, \work{Enarratio in Psalmum} 85 §§1--6 (PL~37, 1081--1090)}
This is the locus classicus of the \latin{totus Christus}, and the whole of the
Introit's reception runs through it. The Word is head of men and men his members,
so that there is one Christ \latin{qui et oret pro nobis, et oret in nobis, et
oretur a nobis. Orat pro nobis, ut sacerdos noster; orat in nobis, ut caput
nostrum, oratur a nobis, ut Deus noster}. \latin{Inclina, Domine, aurem tuam} is
therefore spoken \latin{ex forma servi} and by the whole Body; \latin{clamavi tota
die} is not one day but \latin{omni tempore \dots\ Unus homo usque in finem
saeculi extenditur}; and the ground of the petition is that \latin{Inclinat aurem,
si tu non erigas cervicem: humiliato enim appropinquat; ab exaltato longe
discedit} (§2). At v.~4 he turns the psalm verse into an
exhortation: \latin{Certe recte admonet membra sua ut sursum cor habeant \dots\
cor ut levetur, voluntatem mutat.}
\end{witnessnote}

\begin{witnessnote}
\wlocus{Cassiodorus, \work{Expositio Psalmorum} on Ps.~85 (PL~70)} A second Latin
witness stating the same reading in his own words: \latin{Ex forma servi (sicut
praefati sumus) Christus omnipotenti supplicat Patri \dots\ quod dicit: Inclina,
ostendit quia se ad ipsum extendere non poterat humana conditio.} On \latin{tota
die}, \latin{totius vitae tempus ostenditur, ut per multa tempora annorumque
curricula quasi unius diei continuus clamor esse monstretur}. At v.~4 he reads the
verse of the servant's own part: \latin{Iucundari petit animam suam fons
hilaritatis et origo laetitiae, ab illa scilicet parte qua servus est.} His text
was read on an optical layer of the Migne volume whose recognition is poor, and
his verse numbering runs one lower than the Vulgate's in psalms carrying an
inscription, so his \latin{Vers.~1} at a titled psalm is the Vulgate's v.~2.
\end{witnessnote}

\textbf{Neither commentator's lemma is the antiphon's at every point, and the
difference reaches the citation.} Augustine's v.~1 has no \latin{ad me} and runs
on into \latin{quoniam egenus et inops ego sum}, the clause the cento does not
sing and on which his §§2--3 hang; his v.~3 reads \latin{Miserere mei} where the
missal prints \latin{mihi}; and his v.~4 reads \latin{Iucunda} where the missal
prints \latin{Lætífica}, his §6 commenting on \latin{iucunda}. Cassiodorus reads
\latin{ad me} and \latin{Laetifica} with the missal but \latin{miserere mei}
against it. So the \latin{ad me} the antiphon adds is older than the Missal and
stands in a sixth-century Latin lemma, while \latin{miserere mihi} stands in no
Bible witness and in no patristic lemma opened here: Swete's Greek at Ps.~85:1
has no pronoun answering
\latin{ad me}, and \latin{eleeson me} at v.~3 is a bare accusative that cannot
decide between a Latin genitive and a dative, so \latin{mihi} for \latin{mei} is a
question of Latin construction and not of what the psalm says.

Within the tracked 1962 registry --- a repository derivative and not the printed
book --- Ps.~85:3 opens the Sixteenth Sunday's Introit, which sings v.~5 with it
and takes Ps.~85:1 as its psalm verse, while Ps.~85:2 \latin{Salvum fac servum
tuum} is the Gradual of the Ember Friday in the first week of Lent --- the head of the psalm here and its
continuation a week later; C3 of the synthesis sets the pair beside the two
Offertories' Ps.~39.

\subsection*{Collect \cue{Coll.}}

\latin{Ecclésiam tuam, Dómine, miserátio continuáta mundet et múniat: et quia
sine te non potest salva consístere; tuo semper múnere gubernétur.} The prayer's
hinge is its semicolon, which divides the causal clause from the single petition
it grounds, so that the two verbs of the first limb stand on one side of it and
the one verb of the second on the other. The uncorrected optical layer of the
Benziger 1962 shows that semicolon too, and those of the Pustet 1862, the Venice
1570 and the 1604 \latin{typica} show a comma in its place.

Three things in it are load-bearing for the rest of the Mass. \latin{Miserátio
continuáta} makes mercy an attribute and not a request, and the participle is the
first word of duration in the three orations. \latin{Sine te non potest salva
consístere} denies the Church any standing of her own; the first unit of the
thematic movement takes that denial together with the Postcommunion's. \latin{Tuo semper múnere} makes the
Church's governance a gift; Augustine, expounding the Gradual's own psalm, uses
the same noun to deny that any progress is a merit, as the second unit of the
thematic movement sets out in his own words.

No Father or Doctor comments on this prayer's text in any source opened here. Its
reception is liturgical, and Honorius's \latin{quam ut Deus mundet, muniat, rogat}
is the nearest thing to an exposition of it.

\subsection*{Epistle \cue{Ep.}}

\subsubsection*{Where the lesson is cut from, and what the cuts carry in}

The lesson opens on a conclusion whose premises it does not print: Gal.~5:25
\latin{Si spíritu vívimus, spíritu et ambulémus} draws together the
flesh-and-Spirit antithesis of 5:16--24, and the pericope closes the same figure
at 6:8 in agricultural terms, so the excerpt is bracketed at both ends by an
argument only one end of which is heard stated. Three threads carry in from
outside it. The \latin{invicem} chain runs 5:13, 5:15 and the lesson's own 5:26.
The law thread runs 5:14, 5:18, 5:23 and the lesson's 6:2 \latin{adimplébitis
legem Christi}. And the lesson stops at 6:10, one verse before Paul takes the pen
at 6:11.

\textbf{Both of the lesson's boundaries fall just clear of the cross.} The lesson
begins one verse after 5:24 \latin{Qui autem sunt Christi, carnem suam
crucifixerunt}, and ends four verses before 6:14 \latin{Mihi autem absit
gloriari, nisi in cruce Domini nostri Jesu Christi}. So the Mass raises vainglory
as a problem and does not read Paul's answer to it, and opposes flesh to spirit
without reading the crucifixion of the flesh that grounds the opposition.

\subsubsection*{Four lexical results}

\textbf{The two burdens are one word in Latin and two in Greek.} The apparent
self-contradiction between \latin{Alter altérius ónera portáte} at 6:2 and
\latin{Unusquísque enim onus suum portábit} at 6:5 is an artefact of the Latin:
the Greek reads \emph{baros} at 6:2 and \emph{phortion} at 6:5, and the Vulgate
renders both by one noun. The Douay follows the Latin and says ``burdens'' and
``burden.''

\textbf{The sowing axiom is a commonplace of the Latin Bible} --- 2~Cor.\ 9:6,
Prov.\ 22:8, Jer.\ 12:13, Eccles.\ 11:4 --- and its verse is numbered differently
in the two languages, the sowing clause standing at the end of Greek v.~7 and the
start of Clementine v.~8.

\textbf{One word in the lesson names a church office, and the English does not
carry it.} \latin{Catechizátur} and \latin{catechízat} at 6:6 transliterate the
Greek verb rather than rendering it, and the Douay reads ``instructed'' and
``instructeth,'' so the Latin and English columns part company there.
\latin{Domésticos fidei} at 6:10 is the household metaphor of Eph.\ 2:19,
\latin{domestici Dei}.

\textbf{And the lesson's verb of carrying is the bier-bearers' verb.} The Greek
\emph{bastazo} occurs twenty-seven times in the New Testament, four of them in
Galatians and two of those inside this lesson, and the appointed text prints
\latin{portáte}, \latin{portábit} and \latin{qui portábant} within one Mass. The
Latin does not distinguish them and neither does the English.

\subsubsection*{What the commentators argue}

\begin{witnessnote}
\wlocus{Jerome, \work{Commentariorum in Epistolam ad Galatas} lib.~III, on Gal.
5:26--6:10 (PL~26, cols.\ 423--432)} At 5:26 he confesses that vainglory infects
almsgiving, long prayer, the pallor of fasting, chastity in every state and
martyrdom itself --- \latin{martyrium ipsum, si ideo fiat ut admirationi et laudi
habeamur a fratribus, frustra sanguis effusus est} (col.~424). At 6:1 he observes
that Paul calls the offender \latin{homo} and does not call the corrector one, and
that Paul wrote \latin{ne et tu tenteris} and not \latin{ne tu cadas}: \latin{Vinci
quippe vel vincere, nonnumquam in nostra est potestate: caeterum tentari, in
potestate tentantis est} (col.~426). At 6:2 he grounds the burden on Ps.~37:5 and
closes \latin{Quae Christi lex est? Hoc est mandatum meum, ut diligatis invicem}
(col.~427). \textbf{At 6:5 he reconciles the two burdens by moving one of them to
the tribunal}: in this life we may help each other by prayers and counsel, but
before the judgement seat \latin{non Job, non Daniel, nec Noe rogare posse pro
quoquam, sed unumquemque portare opus suum} (col.~429). At 6:6 he refutes
Marcion's reading of the verse as being about prayer; at 6:7 he ties \latin{Deus
non irridetur} to v.~6 against the disciple who pleads a failed harvest,
\latin{Excusatio verisimilis hominem potest utcumque placare, Deum non potest
fallere} (col.~430); at 6:8 he answers the Encratite Julius Cassianus, who read
sowing in the flesh as a condemnation of marriage; at 6:10 he makes the time of
sowing the present life (col.~432). \textbf{His lemmata agree with the missal at
the points that matter} --- \latin{Deus non irridetur}, \latin{considerans
teipsum} --- and the Vulgate Galatians is his own revision, so his lemma and the
appointed text are one text. The Latin was read on page images with no text
layer and transcribed by eye.
\end{witnessnote}

\begin{witnessnote}
\wlocus{Augustine, \work{Expositio epistulae ad Galatas} §§52--61 (PL~35,
2145--2148)} Correction of a brother is for Augustine the one test of a spiritual
man, and it may be undertaken only after examining one's own conscience:
\latin{Nihil autem sic probat spiritalem virum quam peccati alieni tractatio};
\latin{Numquam itaque alieni peccati obiurgandi suscipiendum est negotium, nisi
cum internis interrogationibus examinantes nostram conscientiam liquido nobis
coram Deo responderimus dilectione nos facere \dots\ Dilige et dic, quod voles}
(§§56--57). At 6:2 the law of Christ is the law of charity, whence \latin{Eadem
igitur Scriptura et idem mandatum, cum bonis terrenis inhiantes premit servos,
Testamentum Vetus; cum in bona aeterna flagrantes erigit liberos, Testamentum
Novum vocatur} (§58). At 6:3--5 the burden nobody can share is the burden of
conscience: \latin{Non ergo laudatores nostri minuunt onera conscientiae nostrae;
atque utinam non etiam accumulent} (§59). At 6:7--10, \latin{Vident enim
seminationem operum suorum, sed messem non vident \dots\ Bene autem seminare, id
est bene operari, facilius est quam in eo perseverare} (§61). §52 distinguishes
the two vices of 5:26, \latin{aemulatio} healed by peace and \latin{invidia} by
meekness. \textbf{Augustine reads a pre-Vulgate Latin Galatians, and his lemmata
differ from the missal's at nearly every appointed verse}: \latin{spiritu et
sectemur} for \latin{spíritu et ambulémus}, the two participles of 5:26 in reverse
order, \latin{in spiritu mansuetudinis intendens teipsum} for \latin{in spíritu
lenitátis, consíderans teípsum}, \latin{et tunc in seipso} for \latin{et sic in
semetípso}, \latin{proprium onus portabit} for \latin{onus suum portábit},
\latin{Deus non subsannatur} for \latin{Deus non irridétur}, and \latin{non
infirmemur; proprio enim tempore metemus infatigabiles} for \latin{non
deficiámus: témpore enim suo metémus, non deficiéntes}. His §54 argument depends
entirely on his own verb \latin{sectemur}.
\end{witnessnote}

\begin{witnessnote}
\wlocus{Chrysostom, \work{Commentary on the Epistle to the Galatians}, chapter~6
(PG~61, 673--680), read in the nineteenth-century English of the Nicene and
Post-Nicene Fathers} Chrysostom does not make the passage a discipline of the
corrector's conscience but an economy of complementary weaknesses inside one
body. On 6:1: ``He says not, in meekness, but, in a spirit of meekness, signifying
thereby \dots\ that to be able to administer correction with mildness is a
spiritual gift,'' and he arraigns ``the malice of the devil rather than the
remissness of the soul.'' At 6:2 the Church is a building whose stones hold
different positions, and the verb is really ``complete'' --- ``make it up all of
you in common, by the things wherein you bear with one another.'' \textbf{At 6:6
he reads a deliberate institution rather than a rule about stipends}: Christ
imposed on the teacher the necessity of requiring aid from his disciples, in order
to repress his spirit. At 6:10 grace ``invites both land and sea to the board of
charity, only it shows a greater care for its own household.'' He is read here in
nineteenth-century English, so his distinction between ``fulfil'' and ``complete''
is his argument about the Greek.
\end{witnessnote}

\begin{witnessnote}
\wlocus{Aquinas, \work{Super Epistolam ad Galatas lectura} cap.~6 lect.~1--2}
Aquinas reads the lesson as an ordered instruction in how superiors treat
inferiors (6:1), equals equals (6:2) and inferiors superiors (6:6).
\latin{Praeoccupatus} means overtaken by surprise and so the more deserving of
pardon; \latin{in aliquo} excludes the habitual sinner, against whom severity is
right; and Paul wrote \latin{instruite} and not \latin{corrigite} because he
speaks of men overtaken. He gives three modes of bearing another's burden, the
third of them \latin{pro poena sibi debita satisfaciendo, orationibus et bonis
operibus} --- \textbf{a development beyond anything in Jerome, Augustine or
Chrysostom on this verse, and his and not the tradition's}. At 6:8 he reads the
asymmetry of the possessives, and at 6:10 \latin{dum tempus habemus, id est in hac
vita, quae est tempus seminandi}. His second exposition of 6:7 is Jerome's
argument almost word for word, so the two are one continuous tradition and not two
independent witnesses.
\end{witnessnote}

\textbf{Jerome and Augustine make Gal.\ 6:1--2 a discipline of the corrector's own
conscience; Chrysostom makes it an economy of complementary weaknesses inside one
body.} On 6:5, Jerome moves
the unshareable burden to the tribunal, Augustine identifies it with conscience,
Chrysostom with self-regard, Aquinas with the rendering of account. Four
expositors, four answers to one verse.

\subsection*{Gradual \cue{Grad.}}

The chant sings the psalm's exordium and stops. What it does not sing runs from
the sabbath inscription at v.~1 through the wicked springing up as grass at
vv.~7--8 to \latin{Justus ut palma florebit; sicut cedrus Libani multiplicabitur}
(v.~13), \latin{Plantati in domo Domini \dots\ florebunt} (v.~14) and
\latin{Adhuc multiplicabuntur in senecta uberi} (v.~15) --- the harvest imagery
that would have answered the Epistle's sowing, and that nobody hears. The
appointed text departs from the Clementine only in pointing.

\textbf{The chant's first word answers the lesson's last \latin{bonum}.} At Gal.\
6:10 the lesson reads \latin{Ergo dum tempus habémus, operémur bonum ad omnes} and
closes immediately on the five words \latin{máxime autem ad domésticos fidei}, the
catchword standing three words before \latin{máxime}; the responsory printed
directly beneath it on p.~396 begins \latin{Bonum est confitéri Dómino}. Two
things hold the catchword down. In the Greek the lesson uses \emph{kalon} at 6:9
and \emph{agathon} at 6:10, so the triple repetition is partly the Latin's and the
chant answers only the second; and the further pairing with the psalm's palm and
cedar is drawn from verses the chant does not sing.

\latin{Misericórdiam tuam et veritátem tuam} is a fixed collocation, standing in
the tracked Clementine at Ps.\ 39:11, Ps.\ 87:12 and Ps.\ 91:3 and nowhere else
--- two of the three in psalms this Mass sings from.

\begin{witnessnote}
\wlocus{Augustine, \work{Enarratio in Psalmum} 91 §§1--4 (PL~37, 1171--1178)} The
sabbath of the title is interior: \latin{Intus est, in corde est sabbatum nostrum
\dots\ Cui autem bona est conscientia, tranquillus est}, and against a merely
bodily observance, \latin{Melius est enim arare, quam saltare}. \textbf{Confessing
is two-sided}, and this is the reading the Gradual's first word carries: \latin{In
utraque re, et in peccato tuo, quia tu fecisti, et in bono facto, confitere
Domino, quia ipse fecit}, with the excuses of Satan and of the stars each refused.
\latin{Psallere nomini tuo} is distinguished from singing --- \latin{Opus nostrum,
psalterium nostrum est \dots\ Canta ore, psalle operibus} --- and the singing seeks
\latin{gloriam Dei, non tuam; nomen ipsius, non tuum}, which is in praise exactly
what the lesson forbids at Gal.\ 5:26. On the versicle: \latin{Mane dicitur,
quando nobis bene est: nox dicitur, quando tristitia tribulationis est \dots\
quando bene est, lauda misericordiam; quando male, lauda veritatem: quia peccata
flagellat, non est iniquus.} His lemma matches the Gradual's wording at both
verses, including \latin{Altissime}, differing only in the comma the 1962 book
drops.
\end{witnessnote}

\begin{witnessnote}
\wlocus{Cassiodorus, \work{Expositio Psalmorum} on Ps.~91 (PL~70)} The same
two-sided confession stated in his own words: confession is of sins and of benefits,
\latin{nec meritorum nostrorum esse dicimus quod ipsius miseratione praestatur};
\latin{Psallere enim significat piis operibus Domini mandata peragere, ut sicut
psalterium de superioribus sonat, ita operatio nostra ad aures Divinitatis
ascendat}. On the versicle: \latin{Mane significat gaudium, nox autem cognoscitur
indicare tristitiam \dots\ Per noctem vero, id est per tribulationem \dots\
veritatem ipsius similiter annuntiare debemus: quia merito patimur quod delictis
gravantibus sustinemus.}
\end{witnessnote}

\begin{witnessnote}
\wlocus{Hilary of Poitiers, \work{Tractatus de titulo psalmi XCI} §§1--3 (CSEL~22,
pp.~339ff.)} \textbf{What survives of Hilary on this psalm is a tractate on its
title, not a commentary on its verses}, and he has nothing on \latin{Bonum est
confitéri Dómino} or on \latin{Ad annuntiándum mane misericórdiam tuam} --- the
two clauses the Gradual actually sings. What he argues is that the sabbath the
title names cannot be the seventh day of Genesis, because the psalm says nothing
about that rest: \latin{nulla hic diei septimi requies, nullum humanorum operum
otium decantatur}. The rest meant must therefore be sought at Ps.\ 94:11, and the
Lord himself loosed the seventh-day sabbath when he said that the priests in the
temple profane it and are blameless. Zingerle's preface records the tractates on
the titles of Pss.\ 9 and 91 as transmitted \latin{solis codicibus Vr}.
\textbf{Hilary's Latin reaches this guide through a poor optical layer of a scan
of Zingerle's edition}, whose running heads and apparatus are badly mangled and
whose page sequence is unreliable.
\end{witnessnote}

Within the tracked registry the same antiphon and verse range stands at the
Saturday of the second week of Lent, and Ps.~91:2 at Septuagesima as an
Offertory; Schuster states the Lenten identity independently at printed p.~139 of
volume~III of \work{The Sacramentary}.

\subsection*{Alleluia \cue{All.}}

Ps.~94 is a summons to worship in vv.~1--7 and an oracle of warning in vv.~8--11,
the turn falling at v.~7/8; the chant's verse stands inside the summons as the
first of its reason clauses, four verses before the turn, so the chant carries in
the summons and not the warning. The warning half is what Hebrews quotes at length, at Heb.\ 3:7--11, 3:15, 4:3,
4:5 and 4:7, three of the five returning to \latin{Hodie si vocem ejus
audieritis}.
The Fourteenth Sunday's Alleluia is \latin{Venite, exsultemus}, Ps.~94:1, so this
chant continues the previous Sunday's psalm --- Schuster calls it ``a
continuation'' and not the next verse, and it is not: 94:1 to 94:3 skips v.~2.

\textbf{No checked witness expounds the predicate this chant substitutes.} What
is proper to this element is that the two expositions of Ps.~94:3 opened here are
built on the word the chant removes, so that neither survives it; and the three
composed orations have no expositor in any source opened here. The chant reads
\latin{et Rex magnus super omnem terram} where every Bible witness opened here
reads \emph{above all gods} ---
Swete's \emph{epi pantas tous theous}, the Clementine's \latin{super omnes deos},
the Douay's ``above all gods.''

\begin{witnessnote}
\wlocus{Augustine, \work{Enarratio in Psalmum} 94 §§5--6 (PL~37, 1217--1234)}
Augustine's lemma is \latin{et rex magnus super omnes deos}, quoted three times in
§§5--6 and again at the close of §6, and he builds two whole developments on the
word the chant removes. First, with 1~Cor.\ 8:5--6 and Ps.\ 95:5, the gods of the
nations are demons, over whom it would be small praise for God to be great:
\latin{Parum enim erat quia terribilis super omnia daemonia Deus: quid magnum,
esse super omnia daemonia?} Second, the \latin{gods} are men made gods by
participation: \latin{Hic accipe homines deos \dots\ Deos dixit participatione, non
natura; gratia, qua voluit facere deos \dots\ Deus verus facit deos credentes in
se.} \textbf{Neither development survives the chant's substitution}, so both are
exposition of Ps.\ 94:3 as the Latin Bible reads it and not of the Alleluia as
sung. His §8 reading of the neighbouring \latin{in manu eius fines terrae}, v.~4,
of Christ joining two peoples is adjacent in sense to the chant's wording and belongs
to a verse the chant does not sing.
\end{witnessnote}

\begin{witnessnote}
\wlocus{Cassiodorus, \work{Expositio Psalmorum} on Ps.~94 (PL~70)} Read after
Augustine and on the same verse, Cassiodorus also has \latin{super omnes deos} and
expounds the gods the same two ways: \latin{Deos frequenter legimus ab hominibus
confictos, ut Iovem, Martem, Saturnum \dots\ Legimus etiam deos, sanctos viros a
Domino constitutos, ut est illud: Deus stetit in synagoga deorum.} So two Latin
commentators expound this verse, and both expound the predicate the chant does
not sing.
\end{witnessnote}

\textbf{The chant's predicate is the Latin Bible's own phrase.} It stands
verbatim at Ps.\ 46:3, \latin{quoniam Dominus excelsus, terribilis, rex magnus
super omnem terram}, the one place in the tracked Clementine where \latin{rex
magnus} and \latin{super omnem terram} collocate.

Two continuous expositions of Ps.\ 94:3 stand here and both expound a different
predicate; no witness located expounds the chant's. In the older chant record the
verse is cued at eight places in the antiphonaries Hesbert
collected, none of them a Sunday after Pentecost later than the Eighth, and the
Ottobonianus margins cue it at a Paschal Friday and at another Mass.
\textbf{Neither phrase the English Bibles carry renders this chant}: ``a great
King above all gods'' is the Bible text the missal departs from, and ``a great
king over all the earth'' is the Douay's Ps.\ 46:3, whose Latin the chant's
predicate reproduces. An identified English of the chant's own reading exists
outside the Bibles: the 1861 Philadelphia hand missal prints the versicle at
printed p.~428 as ``a great King over all the earth,'' beside the Latin
\latin{Quóniam Deus magnus Dóminus, et Rex magnus super omnem terram}. Within
the 1962 temporal registry the chant occurs at this Mass and nowhere else.

\subsection*{Gospel \cue{Gosp.}}

\subsubsection*{Seven results from the text itself}

\textbf{The missal's incipit cuts the connective.} Luke's \latin{Et factum est:
deinceps ibat} makes the raising the next thing in order after the centurion's
servant, and \latin{In illo témpore: Ibat Iesus} replaces it while supplying the
subject the sentence needs. Two consequences follow: the pair Luke makes --- a
servant healed at a distance for a Gentile who asks, then a son raised at close
quarters for a widow who does not --- is severed; and \latin{mortui resurgunt} at
Lk.\ 7:22, six verses after the pericope ends, has no antecedent in the Third
Gospel but this. \latin{Naim} is a hapax in the tracked Clementine.

\textbf{The narrative is told in the words of Elias at Sarepta, and the identity
is exact only in the Greek.} The clause \emph{kai edoken auton te metri autou}
stands at LXX 3~Kingdoms 17:23 and is letter for letter the clause Luke 7:15
ends with; the Greek verse runs on past it into Elias's own words. The Vulgate
breaks the link, reading \latin{tradidit matri suæ} at 3~Kings and \latin{dedit
illum matri suæ} at Luke, so a reader of the Latin column cannot see the
quotation, and the claim that Luke is quoting is a claim about the Greek. Two weaker supports stand beside it: the gate-and-widow agreement at
3~Kingdoms 17:10 against Luke 7:12 uses cognate but different nouns, and
4~Kings~4 shares the raising and the restoring but no verbal formula. The
strongest text-internal support for the Elias figure is Luke's own, at 4:25--26,
where the Sarepta widow is the programmatic example of the Nazareth sermon.

\textbf{\latin{Deus visitávit plebem suam} answers the Benedictus.} The verb and
object repeat Zachary's \latin{quia visitavit, et fecit redemptionem plebis suæ}
at Lk.\ 1:68. \latin{Visitare} with God as its subject stands at Lk.\ 1:68, 1:78,
7:16 and 19:44 \latin{eo quod non cognoveris tempus visitationis tuæ} and nowhere
else in Luke, the Latin stem and the Greek lemmata agreeing verse for verse. The
word enters Luke twice in the
Benedictus, is spoken once by the crowd at Naim, and returns once over Jerusalem.
\textbf{It is in the mouth of the crowd and not of the Evangelist.}

\textbf{\latin{Filius únicus} renders \emph{monogenes}, and the Latin hides the
connection.} \emph{Monogenes} stands nine times in the Greek New Testament ---
Lk.\ 7:12, 8:42, 9:38; Jn.\ 1:14, 1:18, 3:16, 3:18; Heb.\ 11:17; 1~Jn.\ 4:9 ---
and the Latin split is complete with no exception: \latin{unicus}
in Luke, \latin{unigenitus} everywhere else.

\textbf{\latin{Misericórdia motus} renders \emph{splanchnizomai}, which Luke uses
exactly three times}: here, of the Samaritan at 10:33, and of the prodigal's
father at 15:20. The cognate noun occurs once in Luke, at 1:78 \latin{per viscera
misericordiæ Dei nostri} --- the same Benedictus that supplies the visitation
language, and there the Latin does keep the anatomy. So the Naim pericope answers
the Benedictus twice, at v.~13 and at v.~16, and only the second is visible in the
Latin. \latin{Noli flere} and its plural return in Luke exactly three times: 7:13,
8:52 and 23:28.

\textbf{Lk.\ 7:13 is the first place in the Third Gospel where the narrator's own
voice calls Jesus \latin{Dóminus}}, and the appointed text keeps it: before 7:13
the word is God in narration or in cited Scripture, or a character's vocative, or
inside a character's speech.

\textbf{\latin{Prophéta magnus} is not a stock biblical title.} In the tracked
Clementine it stands only at Ecclus.\ 48:25, of Isaias, and at this verse. \latin{Surge} at v.~14 and \latin{surréxit} at v.~16 are one Greek verb in
two forms; the Vulgate's promise of a prophet at Deut.\ 18:15 stands on a
different Latin verb, \latin{suscitabit}.

\subsubsection*{What the commentators argue}

\begin{witnessnote}
\wlocus{Ambrose, \work{Expositio Evangelii secundum Lucam} V.88--92 (CSEL~32
pars~4, pp.~216--218; PL~15, 1376B--1377A)} Three moves at once. \textbf{The widow
is holy Church}, who \latin{populum iuniorem a pompa funeris atque supremis
sepulchri suarum reuocet ad uitam contemplatione lacrimarum, quae flere prohibetur
eum cui resurrectio debeatur}, and ringed by the crowd of peoples she seems
\latin{plus \dots\ esse quam feminam} (§89). \textbf{The bier is wood, and because
it is wood the dead man already has hope of rising}: \latin{spem resurgendi
habebat, quia ferebatur in ligno \dots\ ut esset indicio salutem populo per crucis
patibulum refundendam} (§90). And the bearers who halt are the mortal fluxes of
material nature --- \latin{cum uel ignis inmodicae cupiditatis exaestuat uel
frigidus umor exundat uel pigra quadam corporis habitudine uigor habetatur animorum
\dots\ hi sunt nostri funeris portitores}. At §§91--92, \latin{quis iste est
tumulus tuus nisi mali mores? tumulus tuus perfidia est, sepulchrum tuum guttur
est} (Ps.\ 5:11), and \latin{fleat pro te mater ecclesia, quae pro singulis
tamquam pro unicis filiis uidua mater interuenit}. At VI.64, outside the
pericope's own place, he explains why this raising is public where Jairus's
daughter's is private and reads the two as two peoples: \latin{est etiam forma
sapientiae in uiduae filio cito ecclesiam credituram, in archisynagogae filia
credituros quidem Iudaeos, sed ex pluribus pauciores} --- marked with \latin{puto}
and \latin{est etiam forma}, offered as a reading and not as a demonstration. The
critical text was read as an optical layer whose body text is generally sound and
whose apparatus lines are mangled.
\end{witnessnote}

\begin{witnessnote}
\wlocus{Augustine, \work{Sermo} 98 (PL~38, 591--595)} \textbf{The three raisings
are three grades of sin.} The girl inside the house is sin consented to in the
heart; the young man carried out of the gate but not yet buried is sin gone out
into the deed; Lazarus, buried under the stone, is sin hardened into habit ---
\latin{Moles illa imposita sepulcro, ipsa est vis dura consuetudinis, qua premitur
anima, nec surgere, nec respirare permittitur} (§5). The four-step analysis behind
it stands at §6: \latin{Prima est enim quasi titillatio delectationis in corde;
secunda, consensio; tertium, factum; quarta, consuetudo.} At §2 the widowed mother
is mother Church rejoicing at each raising; at §6 the loosing of Lazarus is the
office given to the disciples while the raising stays Christ's. \textbf{The sermon
stops at v.~15}, one verse short of the appointed pericope, and nothing in it
comments on v.~16.
\end{witnessnote}

\begin{witnessnote}
\wlocus{Cyril of Alexandria, \work{Commentary on the Gospel of Saint Luke}, Sermon
XXXVI (Payne Smith, Oxford 1859, vol.~1, pp.~132--135)} Christ comes uninvited:
at Capharnaum he was present by invitation, but at Naim ``He draws near without
being invited. For no one summoned Him to restore the dead man to life, but He
comes to do so of His own accord.'' The Evangelist supplies two proofs of a real
raising, that the dead man began to speak and that Christ gave him to his mother;
the three raised are ``a pledge of the hope prepared for us of a resurrection of
the dead,'' confirmed from Is.\ 26:19 and Ps.\ 103(104):29--30. \textbf{And the
touching of the bier teaches the life-giving flesh}: it was done ``that thou
mightest learn that the holy body of Christ is effectual for the salvation of man.
For the flesh of the Almighty Word is the body of life \dots\ For consider, that
iron, when brought into contact with fire, produces the effects of fire,'' so that it
``annihilates the influence of death and corruption.'' The words reach this page
in an English translation of an ancient Syriac version of a Greek original that
survives elsewhere only in catena fragments.
\end{witnessnote}

\begin{witnessnote}
\wlocus{Bede, \work{In Lucae evangelium expositio} lib.~II (PL~92, cols.\
417--420)} Bede takes Augustine's moral reading and works it through the
pericope's furniture. At v.~12 the dead man carried out is sin no longer hidden in
the heart's chamber but published abroad, \latin{quasi per suae civitatis ostia
propalantem}, and the widowed mother is mother Church, one perfect and immaculate
virgin though gathered from many persons, \latin{licet e multis collecta personis,
una sit perfecta et immaculata virgo, mater Ecclesia}, each of the faithful her
only son, with Gal.\ 4:19 quoted for her motherhood; the gate is one of the bodily
senses, closing with \latin{Obsecro, Domine
Jesu, cunctas meae civitatis portas justitiae facias} (Ps.\ 117:19). Against
Novatian, who \latin{veramque matrem Ecclesiam \dots\ negat consolari debere}. At
v.~14 the bier is \latin{male secura desperati peccatoris conscientia} and its
bearers \latin{vel immunda desideria \dots\ vel lenocinia blandientium \dots\
venenata sociorum}; at v.~15 \latin{Redditur matri, cum \dots\ per sacerdotalis
decreta judicii communioni sociatur Ecclesiae}. \textbf{Bede is the only witness
located who comments on v.~16}: \latin{Quia Deus visitavit plebem suam. Non
tantummodo verbum suum semel incorporando, sed etiam \dots\ semper in corda
mittendo.} Two of his moves are his own and not Ambrose's or Augustine's: the
anti-Novatianist application, and the reading of the restoration to the mother as
sacramental reconciliation. His text was read on an optical layer of Migne whose
two-column setting interleaves badly and whose individual words are corrupt in
places.
\end{witnessnote}

\begin{witnessnote}
\wlocus{Gregory of Nyssa, \work{De hominis opificio} 25.6--11 (PG~44, 213--217),
read in the nineteenth-century English of the Nicene and Post-Nicene Fathers,
second series vol.~5} The raising is fourth in a graded pedagogy by which Christ
accustomed human weakness to believe in resurrection --- the fever,
the nobleman's son, the ruler's daughter, the young man on his way to the tomb,
four-day Lazarus --- and the widow's only son is ``the stock of her race, the shoot
of its succession, the staff of her old age.''
\end{witnessnote}

\textbf{Four expositors, four questions.} Augustine reads the raisings morally,
Cyril eschatologically, Ambrose ecclesially, Gregory of Nyssa as a pedagogy of
increasing difficulty. None denies another and they answer different questions.
On the widow the two Latin readings are not identical
either: Ambrose's Church weeps and intercedes for each sinner as for an only son
and is forbidden to weep because resurrection is owed to her; Augustine's
rejoices. Grief and joy, not one doctrine.

\subsubsection*{What the commentators do not supply}

Gregory the Great has nothing on this pericope. Within Luke 7 every heading of
his in the \work{Catena aurea}, bare \latin{Gregorius} or \latin{Gregorius in
Evang.}, stands at lect.~6, on the sinful woman of Lk.\ 7:36--50, and none at
lect.~2; the \latin{Gregorius Nyssenus} headings are another Gregory's. On the
standard account of their contents, no edition of
them having been opened for this formulary, Gregory's forty \work{Homiliae in
Evangelia} treat Luke~7 only at Homily~33, on the same pericope of the sinful
woman.

\textbf{The eight-resurrections scheme the \work{Catena aurea} attributes to
\latin{Ambrosius} is not in Ambrose's Luke commentary}: a word search of the
complete critical text for the scheme's own terms returns V.86 and VI.64 only,
and neither carries it; it may stand in another work of his, or be Ps.-Ambrose,
or be a Glossa attribution. Three further attributions in the same catena at this chapter have no
located text behind them --- Theophylact's widow-as-soul allegory, Titus of
Bostra's contrast with Elias and Eliseus, and the unnamed Chrysostom fragment on
consolation.

\correction{Anthony of Padua preached this Gospel at his Sixteenth Sunday, not his Fifteenth}{%
Anthony's sermon on the widow of Naim is his \latin{Dominica XVI post
Pentecosten}, which opens \latin{Evangelium in decima sexta dominica post
Pentecosten: Ibat Iesus in civitatem quae vocatur Naim} and pairs the pericope
with Eph.\ 3:13--21; his \latin{Dominica XV} opens \latin{Nemo potest duobus
dominis servire}. A citation made by matching the ordinal would name the wrong
sermon silently.}

The same pericope is appointed at the Thursday of the fourth week of Lent in the
tracked registry, where the Missal pairs it with 4~Kings 4:25--38, Eliseus and the
Sunamitess's son; Schuster and Guéranger both name that Lenten Thursday
independently, and the Pustet 1862 does not print this Sunday's Gospel at all but
cross-references it there.

\subsection*{Offertory \cue{Off.}}

The antiphon takes vv.~2, 3 and 4 of Ps.~39 and cuts the middle out of them: the
whole of v.~3 from \latin{et eduxit me de lacu miseriæ} to \latin{direxit gressus
meos} goes, and so does v.~4b. What lies immediately past the second cut is
\latin{Videbunt multi, et timebunt, et sperabunt in Domino} and, one verse on,
\latin{Beatus vir cujus est nomen Domini spes ejus}. Two verses further the
psalm's vv.~7--9 are what Hebrews puts in the mouth of Christ coming into the
world, and there the Latin Bible disagrees with itself across the seam: the psalm
reads \latin{aures autem perfecisti mihi} where Heb.\ 10:5 quoting it reads
\latin{corpus autem aptasti mihi}, while Swete's Greek Psalter reads \emph{soma}
with Hebrews and against the Latin psalm. The antiphon stops three verses short of
that seam, at \latin{Deo nostro}, so nothing in its own extent rests its place at
the offering of the gifts on Hebrews.

\textbf{Three of the antiphon's phrases are not the Clementine's, and all three
have Latin witnesses no later than the sixth century.} Against Swete, the second and third are the
literal renderings and the Clementine's are the freer: \emph{deeseos} is a
singular that \latin{deprecationem meam} renders exactly where \latin{preces meas}
does not, and \emph{hymnon} is \latin{hymnum} and not \latin{carmen}. The first,
\latin{respéxit me}, answers neither Swete's \emph{proseschen moi} nor any Greek
variant collated here --- but it stands in Cassiodorus's lemma, so it has a witness
even though it has no derivation.

\begin{witnessnote}
\wlocus{Augustine, \work{Enarratio in Psalmum} 39 §§2--5 (PL~36, 433--450)} The
dative is Augustine's point, and his lemma is not the antiphon's at either of
its first two verbs, reading \latin{Sustinens sustinui} for \latin{Exspéctans
exspectávi} and \latin{attendit mihi} for \latin{respéxit me}: \latin{Ergo et
cum malus esses, attendebat te, sed non attendebat tibi \dots\ Attendit mihi,
inquit, id est, consolando attendit, ut mihi prodesset} (§2). The pit of misery
--- which the antiphon cuts --- is \latin{Profunditas iniquitatis, ex carnalibus
concupiscentiis}, and the rock is Christ (§3). The new song is the song of the new
man: \latin{novus homo, dicat canticum novum: innovatus amet nova quibus renovatur
\dots\ Dicimus ergo hymnum Deo nostro, et ipse hymnus liberat nos} (§4) ---
\textbf{his lemma reading \latin{hymnum} with the antiphon and against the
Clementine}. At §5 he answers the question of the psalm's speaker outright: it is
Christ, speaking now in the name of our Head and now of us who are his members.
The exposition opens with a long anti-Donatist prologue at §1 that has nothing to
do with the appointed verses.
\end{witnessnote}

\begin{witnessnote}
\wlocus{Ambrose, \work{In Psalmum XXXIX enarratio} §§1--4 (PL~14, cols.\
1109--1110)} \textbf{The doubled verb is the point}: the psalm announces the New
Testament, and \latin{non est perfecti exspectare, sed exspectasse; nemo enim nisi
qui perseveraverit usque in finem, salvus erit} (§1) --- the same Matt.\ 10:22
that Jerome and Augustine both reach for at the lesson's \latin{non deficiámus}.
At §3, \latin{immisisti in os meum canticum novum, quod est Novum Testamentum.
Gaudentes iam canimus hymnum Deo nostro, quia novarum virtutum praecepta
cognovimus; ut nostra omnia relinquamus, Christum sequentes, et nostros diligamus
inimicos \dots\ Maledicentes ecce benedicimus.} His lemma carries the antiphon's
incipit and \latin{hymnum}. The Latin was read on page images with no text layer
and transcribed by eye.
\end{witnessnote}

\begin{witnessnote}
\wlocus{Cassiodorus, \work{Expositio Psalmorum} on Ps.~39 (PL~70)}
\textbf{Cassiodorus's lemma is the antiphon's text word for word} ---
\latin{Exspectans exspectavi Dominum, et respexit me: et exaudivit deprecationem
meam} --- and he reads \latin{hymnum} for the Clementine's \latin{carmen}. The
speaker is the Church: \latin{Ecclesia catholica, quae fuerat de totius mundi
partibus congreganda, patientiae virtutem praedicat.} On the doubling,
\latin{considerandus est hic sermo geminatus, quia superflua non est tam decora
repetitio}; on the verb the Clementine does not have, \latin{Respexit: quia
praevenit omne bonum nostrum}; on the noun, \latin{Deprecatio quoque significat
frequentissimam precem}; and on \latin{hymnum Deo nostro}, \latin{Hymnus autem
Graecus sermo est, id est, laus carminum lege composita. Et quoniam hymni erant
quos idolis suis etiam gentilitas personabat, addidit, Deo nostro.}
\end{witnessnote}

\textbf{Ps.~39 and Ps.~69 stand at three Sundays inside a five-Sunday span.} The
Twelfth Sunday's Introit sings Ps.\ 69:2--3, this Offertory Ps.\ 39:2--4 and the
Sixteenth Sunday's Offertory Ps.\ 39:14--15, one psalm supplying two consecutive
Offertories. The doublet between the two psalms, and with it Cassiodorus's
division of the repeated words between speakers --- Christ's in his own person at
Ps.~39 and his members' at Ps.~69, which neither psalm's own text marks --- stand
at their loci in the guide to the Twelfth Sunday, whose Introit sings that psalm,
and were not opened again here: Cassiodorus was read for this Mass at Pss.\ 39,
85, 91 and 94. Within the tracked registry the antiphon also stands at the
Tuesday of the fourth week of Lent, which Schuster states independently at printed
p.~140.

\subsection*{Secret \cue{Sec.}}

\latin{Tua nos, Dómine, sacraménta custódiant: et contra diabólicos semper
tueántur incúrsus.} Two verbs of defence, one object, and an adverb that has
moved: both ancient Roman witnesses read \latin{tueantur semper incursus} where
1962 reads \latin{semper tueántur incúrsus}. The prayer's object is the
congregation and no longer the Church, one step of the narrowing the fourth unit
of the thematic movement reads across the three orations.

\textbf{No exposition of this prayer's text stands in any source opened for this
formulary.} What the prayer asks to be defended from is named in
the lesson it follows: Chrysostom, on Gal.\ 6:1, arraigns ``the malice of the
devil rather than the remissness of the soul'' in the man overtaken in a fault,
and that malice is the \latin{diabólicos \dots\ incúrsus} of the Secret. He is
commenting on Paul.

\subsection*{Communion \cue{Comm.}}

The antiphon sings the exact clause at which the Bread of Life discourse turns
from bread to flesh, and only its second half: the half not sung is \latin{Si quis
manducaverit ex hoc pane, vivet in æternum}, and the words that are sung are the
ones that make the crowd start quarrelling in the next verse, \latin{Litigabant
ergo Judæi ad invicem, dicentes: Quomodo potest hic nobis carnem suam dare ad
manducandum?} The element the excerpt carries in and cannot state is the
discourse's resurrection refrain: \latin{et ego resuscitabo eum in novissimo die}
three times, at vv.~40, 44 and 55 in the missal's own Clementine numbering, and
the same promise in other words at v.~39, \latin{sed resuscitem illud in
novissimo die}.

\textbf{Both of the antiphon's departures are on the Latin side.} The Greek has a
simple future \emph{doso} and \emph{kosmou}, so \latin{dédero} and \latin{sǽculi}
answer nothing in it. Neither is an error --- \latin{dedero} is a future
perfect and both \latin{sæculum} and \latin{mundus} are live Latin renderings of
the Greek noun elsewhere in the Vulgate.

\begin{witnessnote}
\wlocus{Augustine, \work{In Iohannis evangelium tractatus} 26 §§12--18 (PL~35,
1607--1616)} The tractate's own printed heading is the antiphon's verse in the
Clementine's wording, \latin{Panis quem ego dabo, caro mea est pro mundi vita}.
The eating that gives life is inward: \latin{qui manducat intus, non foris; qui
manducat in corde, non qui premit dente} (§12). Only the body of Christ lives by
the Spirit of Christ: \latin{De Spiritu Christi non vivit, nisi corpus Christi
\dots\ O Sacramentum pietatis! o signum unitatis! o vinculum caritatis! Qui vult
vivere, habet ubi vivat, habet unde vivat. Accedat, credat ut incorporetur, ut
vivificetur} (§13). The \latin{res} of the sacrament is \latin{societatem corporis
et membrorum suorum, quod est sancta Ecclesia}, and it is \latin{omni homini ad
vitam, nulli ad exitium} where the sacrament is \latin{quibusdam ad vitam,
quibusdam ad exitium} (§15); the species is made of many grains and many berries
(§17). \textbf{His lemma is \latin{dabo} and \latin{pro mundi vita}}, so this is
exposition of the verse and not of the antiphon as sung.
\end{witnessnote}

\begin{witnessnote}
\wlocus{Cyril of Alexandria, \work{Commentary on the Gospel according to S. John}
IV.2 (Pusey, Library of the Fathers 43, Oxford 1874, vol.~1, pp.~409--412)} The
chapter heading is the doctrine: \emph{That the Holy Body of Christ is
Life-giving}. The indwelling Word ``transformed it into His Own proper good, that
is life \dots\ rendered It life-giving, as Himself is by Nature,'' with Jn.\ 17:19
read alongside and the heave-offering of Num.\ 15:18--21 as its figure. This is
the same doctrine Cyril states at the Gospel's bier, argued there from the
touching and here from the Eucharist. The words are Pusey's English of the
Greek.
\end{witnessnote}

\begin{witnessnote}
\wlocus{Chrysostom, \work{Homily} 46 \work{on the Gospel of John} (PG~59,
257--260), read in nineteenth-century English} ``By bread He means here either His
saving doctrines and the faith which is in Him, or His own Body; for both nerve
the soul.'' On the appointed clause: ``observe how little by little He led them up
to Himself. Here He says that Himself gives, not the Father.'' And on the
quarrelling that follows it, ``when questioning concerning the how comes in, there
comes in with it unbelief.''
\end{witnessnote}

\begin{witnessnote}
\wlocus{Aquinas, \work{Super Evangelium S. Ioannis lectura} cap.~6 lect.~6} The
flesh is \latin{organum divinitatis suae}, and the instrument acts by the power of
the agent: \latin{cum instrumentum agat virtute agentis \dots\ ita ut Damascenus
dicit et caro virtute verbi adiuncti vivificat: unde Christus tactu suo sanabat
infirmos}. \textbf{That last clause reaches the Gospel of this same Mass}, where
Christ touches the bier. The bread is the sacrament of the Church's unity because
made of many grains; on the verb, \latin{dicit dabo, quia nondum institutum erat
hoc sacramentum}; \latin{Non dicit autem carnem meam significat sed caro mea est};
and on \latin{pro mundi vita}, \latin{vita quam confert \dots\ quantum in se est,
totius mundi}. \textbf{His gloss on the tense turns on the simple future and
cannot be transferred to \latin{dédero}}, nor his exposition of the universality
of the effect to \latin{sǽculi}.
\end{witnessnote}

\correction{Two numbering series meet at this antiphon}{%
The missal cites the antiphon as Io.\ 6, 52, and the tracked Clementine and Douay
both number it 52, so missal and Bible agree without adjustment. The Greek Fathers
and the nineteenth-century English translations made from them number this clause
6:51 and reserve 6:52 for \latin{Quomodo potest hic nobis carnem suam dare ad
manducandum?} A citation of Cyril's chapter by its modern label would therefore be
wrong by one verse against the missal's reference unless the series is named.}

Within the tracked registry the antiphon also stands at the Thursday of the first
week of Lent; Schuster states the same at printed p.~141 and adds that the Ninth
Sunday's antiphon comes from the same discourse, though his reference there is
loose, the Ninth Sunday's being \latin{Qui manducat meam carnem} at Jn.\ 6:57.

\subsection*{Postcommunion \cue{Postcomm.}}

\latin{Mentes nostras et córpora possídeat, quǽsumus, Dómine, doni cæléstis
operátio: ut non noster sensus in nobis, sed iúgiter eius prævéniat efféctus.}
\textbf{The subject is held back to the end of its clause}, behind two objects and
a vocative: it is \latin{doni cæléstis operátio} that is asked to possess, and its
\latin{efféctus} --- not \latin{noster sensus} --- that is asked to go before. The
pointing carries the grammar: no comma after \latin{nostras}, so \latin{Mentes
nostras et córpora} is a single object pair, and the colon after \latin{operátio}
divides the petition from the purpose clause.

\subsection*{The three orations, and their liturgical history}

No Father or Doctor comments on \latin{Ecclésiam tuam, Dómine, miserátio
continuáta}, \latin{Tua nos, Dómine, sacraménta custódiant} or \latin{Mentes
nostras et córpora possídeat} in any tracked source, or in any patristic or
medieval work opened for this formulary.

\textbf{Their history is liturgical, and it begins where the books begin.} None of
the three occurs in the Veronense, in Férotin's \work{Liber mozarabicus
sacramentorum}, in Bannister's \work{Missale Gothicum} or in Lowe's \work{Bobbio
Missal}. \latin{Incursus} and \latin{ecclesiam tuam} stand in both Veronense
scans and in the Mozarabic and the Gothicum, \latin{possideat} and \latin{mentes
nostras} in three of those four layers, and \latin{continuata} in the Veronense
scans alone; no control count stands for the Bobbio Missal, and an optical
layer that renders the \latin{e caudata} as a caret would hide any of the three
prayers in any of the four books. No source opened for this formulary dates or
attributes any of the three: their attested history begins at the Old Gelasian and
the Hadrianum.
